\clearpage
\makeatletter
\let\@oldaddcontentsline\addcontentsline
\renewcommand{\addcontentsline}[3]{}

\section*{附录}
\let\addcontentsline\@oldaddcontentsline
\makeatother

\renewcommand{\arraystretch}{1.2}

\subsection*{附录 A：数据利用清单}

附录 A 按附件编号说明每份数据在本文中的用途，并把真实记录、半合成实验、插值轨迹和外推估算区分开来。这样既方便复核文件与结果的对应关系，也能看出哪些结论来自直接观测，哪些只用于对照或稳健性检查，详见表\ref{tab:appendix}。

\begin{longtable}{@{}>{\centering\arraybackslash}p{0.09\textwidth}>{\raggedright\arraybackslash}p{0.23\textwidth}>{\centering\arraybackslash}p{0.08\textwidth}>{\centering\arraybackslash}p{0.08\textwidth}>{\raggedright\arraybackslash}p{0.40\textwidth}@{}}
  \caption{附件数据利用清单}
  \label{tab:appendix} \\
  \toprule
  编号 & 文件名 & 性质 & 是否使用 & 使用方式与可信度说明 \\
  \midrule
  \endfirsthead
  \caption{附件数据利用清单（续）} \\
  \toprule
  编号 & 文件名 & 性质 & 是否使用 & 使用方式与可信度说明 \\
  \midrule
  \endhead
  \bottomrule
  \endfoot
  A1 & slimpajama\_\allowbreak{}quality\_\allowbreak{}signal\_\allowbreak{}sample & 真实 & 使用 & 对全量 51{,}230 条计算 22 指标熵权评分、领域 $Q$ 与冲突，并抽查原文 \\
  A2 & slimpajama\_\allowbreak{}quality\_\allowbreak{}extended/arxiv & 真实 & 使用 & 以全部 17{,}523 条计算 arxiv 的领域 $Q$，与抽样结果比较 \\
  A3 & slimpajama\_\allowbreak{}quality\_\allowbreak{}extended/github & 真实 & 使用 & 以全部 203{,}752 条计算 github 的领域 $Q$，与抽样结果比较 \\
  A4 & train\_\allowbreak{}mixture\_\allowbreak{}1m & 真实 & 使用 & 512 组训练配方作为岭回归输入，用于拟合配比与损失 \\
  A5 & train\_\allowbreak{}pile\_\allowbreak{}loss\_\allowbreak{}1m & 真实 & 使用 & 提供回归对应的 13 个验证域损失 \\
  A6--A9 & test\_\allowbreak{}mixture/loss\_\allowbreak{}1m, 60m & 真实 & 使用 & 分别检验同尺度和跨尺度的预测，每组 256 条 \\
  A10--A11 & test\_\allowbreak{}mixture/loss\_\allowbreak{}1B & 真实 & 使用 & 用 64 组记录检验跨尺度的预测 \\
  A12--A15 & est\_\allowbreak{}mixture/loss\_\allowbreak{}10b, 70b & 子集/外推 & 使用 & 各 63 组用于比较外推尺度因子的稳定性；属于外推记录，并非观测 \\
  A16 & domain\_\allowbreak{}mapping\_\allowbreak{}guide & 参考 & 使用 & 连接 17 域与质量分类：6 域按 direct/near\_\allowbreak{}direct 对应，11 域标为 inferred 并使用全局均值 \\
  A17 & regmix\_\allowbreak{}domain\_\allowbreak{}summary & 真实 & 未使用 & 未用于本文分析 \\
  A18 & regmix\_\allowbreak{}domain\_\allowbreak{}sample & 真实 & 未使用 & 未用于本文分析 \\
  B1 & pythia\_\allowbreak{}training\_\allowbreak{}log\_\allowbreak{}existing & 真实 & 使用 & 1{,}176 点用于经典拟合和模型留出检验；这组记录的拟合残差极小 \\
  B2 & cerebras\_\allowbreak{}training\_\allowbreak{}log & 半合成 & 使用 & 作模型族外对照，MAPE 为 33.7\%；来源差异较大，只作定性参考 \\
  B3 & training\_\allowbreak{}trajectories & 插值 & 使用 & 检查轨迹内部的 $D$ 指数：$-0.2806$ 至 $-0.2803$，对照拟合值 $-0.2799$ \\
  B4 & scaling\_\allowbreak{}baseline & 真实 & 使用 & 用不同模型族作检验，相关 $r=0.975$，MAPE 为 8.8\% \\
  B5 & published\_\allowbreak{}scaling\_\allowbreak{}data & 真实 & 使用 & 与文献记录比较，相关 $r=0.910$，MAPE 为 8.0\% \\
  B6 & supplementary\_\allowbreak{}NQ\_\allowbreak{}experiment & 半合成 & 使用 & 360 点用于估计质量缺口项；说明半合成条件下的质量效应，不作为直接实验观测 \\
  B7 & supplementary\_\allowbreak{}NQ\_\allowbreak{}experiment\_\allowbreak{}expanded & 半合成 & 使用 & 新增记录用于开发阶段检查，相关 $r=0.991$ \\
  B8 & supplementary\_\allowbreak{}NQ\_\allowbreak{}experiment\_\allowbreak{}large & 半合成 & 部分使用 & $Q_{\text{score}}$ 与 B6/B7 的方向相反，且含低于不可约损失下限的记录；不并入拟合，用于说明来源间的差异 \\
  B9 & supplementary\_\allowbreak{}large\_\allowbreak{}models & 真实 & 使用 & 提供百亿参数以上模型的参数量、数据量和开源可及性资料 \\
  B10 & supplementary\_\allowbreak{}large\_\allowbreak{}baseline & 估算 & 使用 & 比较百亿参数以上的损失估算，$r\approx1.000$、MAPE 为 1.0\%；该对照不作为直接观测 \\
  B11 & open\_\allowbreak{}model\_\allowbreak{}family\_\allowbreak{}metadata & 真实 & 未使用 & 未用于本文分析 \\
  B12 & pythia\_\allowbreak{}checkpoint\_\allowbreak{}index & 真实 & 未使用 & 未用于本文分析 \\
  C1 & leaderboard\_\allowbreak{}cleaned & 真实 & 使用 & 按六维得分、许可、类型和提交日期整理样本；筛选后的 2{,}346 条用于面板回归和前沿分析 \\
  C2 & leaderboard\_\allowbreak{}enhanced & 真实 & 未使用 & 作为 C1 的增强替代版本，为避免重复计入而不另用 \\
  C3 & leaderboard\_\allowbreak{}extended\_\allowbreak{}timeseries & 混合 & 使用 & 核算 2019--2025 年增量；排除 13 条一致性异常记录，长期解释仍保留评分来源差异 \\
  C4 & epoch\_\allowbreak{}all\_\allowbreak{}ai\_\allowbreak{}models & 真实 & 使用 & 观察训练算力趋势，年斜率 $+0.589$ dex；核对数据量与权重字段，Yes 为 1{,}325、No 为 1{,}328 \\
  C5 & loss\_\allowbreak{}benchmark\_\allowbreak{}bridge & 混合 & 使用 & 作为损失与得分关系的对照，$n=43$、$r=-0.502$ \\
  C6 & loss\_\allowbreak{}benchmark\_\allowbreak{}bridge\_\allowbreak{}expanded & 混合 & 使用 & 用于主要损失--得分拟合，$n=75$；High、Medium 两层按可比程度分别处理 \\
  C7 & model\_\allowbreak{}architecture\_\allowbreak{}metadata & 真实 & 使用 & 将 2048、4096、8192、32768、131072 五档上下文用于问题三灵敏度计算 \\
  C8 & detailed\_\allowbreak{}results/ & 真实 & 使用 & 1{,}860 个目录取得有效结果，主样本交集的 964 个用于汇总 37 项任务；六维指标另与完整 C1 逐记录校验 \\
  C9 & data/（Parquet） & 真实 & 未使用 & 与 C1 等价，本文不重复计入 \\
  C10 & pythia\_\allowbreak{}eval\_\allowbreak{}details/ & 说明 & 未使用 & 只有 README 说明，不能作为 C8 的任务记录使用 \\
\end{longtable}

\subsection*{附录 B：半合成与外推数据的标注}

以下资料不全是直接观测，解释相应结果时需分别考虑其生成方式和评分来源。

\begin{itemize}[itemindent=2em]
\item \textbf{半合成数据}：B2 为 Cerebras 训练日志，B6/B7/B8 为 NQ 质量实验。本文以 B6 估计质量项，以 B7 新增点作开发阶段验证，5.2 节表格标明其半合成性质。检查来源差异后，B2、B8 均不参加主拟合。
\item \textbf{外推与估算数据}：A12--A15 的 10B/70B 损失由 1M/60M/1B 记录作幂律外推，B10 则为百亿参数以上的损失估算。它们用来比较估算关系的稳定性，不当作真实训练观测。
\item \textbf{插值数据}：B3 来自 Pythia 检查点之间的插值，只检查轨迹内部幂指数是否一致，不增加独立观测。
\item \textbf{混合口径数据}：C3 各年的评测标准并不完全相同，异常记录已检查并剔除。C5/C6 的 Loss 来源和验证集也未完全统一，因此按可比程度分层使用。
\end{itemize}

\subsection*{附录 C：结果文件与核心求解代码}

本文把求解结果保存为 CSV，位置在 求解/问题一至问题四/结果/ 目录。表\ref{tab:appendix2}列出主要文件及其作用，表中省略“求解/问题X/结果/”这一共同前缀，便于阅读。

\begin{longtable}{@{}>{\raggedright\arraybackslash}p{0.42\textwidth}>{\raggedright\arraybackslash}p{0.54\textwidth}@{}}
  \caption{主要结果文件说明}
  \label{tab:appendix2} \\
  \toprule
  文件 & 说明 \\
  \midrule
  \endfirsthead
  \caption{主要结果文件说明（续）} \\
  \toprule
  文件 & 说明 \\
  \midrule
  \endhead
  \bottomrule
  \endfoot
  质量指标熵权.csv & 22 指标的方向判定、语义说明、信息熵与熵权 \\
  质量信号缺失审计.csv & 逐指标字段级缺失计数（A1--A3 全量） \\
  质量冲突与一致性.csv & 抽样集/扩展集/全量的冲突率、独立基准与 Kendall $W$ \\
  冲突阈值灵敏度.csv & 冲突率随阈值 $\tau$ 的变化 \\
  抽样集与扩展集域级对照.csv & A1 与 A2/A3 的域级质量对照 \\
  质量与损失难度一致性.csv & 两套质量口径在 6 个映射域上的排序比较 \\
  质量项增量检验.csv & 追加混合质量指数后的 $R^2$ 增量 \\
  混合系数矩阵.csv、截距.csv & 配比—损失岭回归系数与截距 \\
  推荐配比调整.csv & 有界推荐配比与调整量 \\
  问题一\_\allowbreak{}拟合汇总.csv & 各损失域的拟合与检验指标 \\
  岭正则敏感性.csv & 正则系数灵敏度 \\
  经典标度律参数.csv & 经典标度律参数、拟合与留出验证指标 \\
  质量字段方向审计.csv & B6/B7/B8 的组内相关与方向结论 \\
  广义标度律形式对比.csv & 四种质量项形式的 $R^2$ 与残差信息评分 比较 \\
  广义标度律分层验证.csv & 6 组数据来源的验证结果 \\
  弹性与替代关系.csv、质量参数替代关系.csv & 弹性与替代率 \\
  领域替代互补\_\allowbreak{}top.csv & 领域效应的相似性与差异 \\
  B3轨迹验证.csv & 轨迹内 $D$ 幂指数验证 \\
  三档预算最优配置.csv、成本函数对比.csv & 三档预算与三型成本函数的优化结果 \\
  结构性转移识别.csv & 结构性转移门槛 \\
  上下文长度敏感性.csv & C7 五档上下文长度下的最优配置 \\
  质量基线灵敏度.csv & 质量基线口径灵敏度 \\
  配比项算力等价.csv & 配比优化的算力等价效果 \\
  筛选漏斗.csv & 字段完整性与许可筛选过程，以及模型类型构成 \\
  C8逐任务解析结果.csv & C8 成功解析目录的任务级原始指标 \\
  C8逐任务聚合.csv、C8子任务类别汇总.csv & 子任务前沿与剩余空间 \\
  C8与C1一致性校验.csv & C8 与 C1 的秩相关校验 \\
  Loss\_\allowbreak{}Benchmark映射.csv & 分层桥接拟合结果 \\
  规模时间分解.csv、规模时间分解\_\allowbreak{}长周期.csv & 规模/非规模贡献分解 \\
  C3口径异常行.csv & C3 自洽性审计剔除的记录 \\
  前沿预测.csv、前沿预测\_\allowbreak{}结构外推情景.csv & 能力前沿预测与情景 \\
  C4宏观趋势.csv、C4宏观摘要.csv & C4 宏观算力与开源权重证据 \\
  求解/广义标度律参数.csv、问题一\_\allowbreak{}关键量.json & 跨问共享的接口参数 \\
\end{longtable}


\subsection*{附录 D：运行环境与核心实现}

运行输入、参数及输出以随文提供的源程序和结果文件为准。各问程序开头注明辅助工具及使用范围。

程序使用 Python 3，主要依赖 NumPy、pandas、SciPy 和 Matplotlib，数据解压及 JSON 读取调用 Python 标准库。四问共用 求解/\allowbreak{}\_common.py，附件路径以该文件设置的数据根目录为准。下面列出四问的完整主程序，涵盖数据读取、模型求解、检验和结果输出；共享函数仍放在求解目录的 \_common.py 中，代码行号与源文件一致。

运行时依次执行问题一、问题二和问题三，问题四读取附件 C、可单独运行。问题一输出质量、配比和关键量，问题二据此估计配比修正及广义标度律参数，问题三读取前两问共享文件。各程序的输入、输出关系见表\ref{tab:code_map}。

\begin{table}[!htbp]
\centering
\caption{问题、程序与数据依赖}\label{tab:code_map}
\begin{tabularx}{\textwidth}{@{}p{0.09\textwidth}>{\raggedright\arraybackslash}p{0.25\textwidth}XX@{}}
\toprule
问题 & 核心程序 & 输入 & 主要输出\\
\midrule
一 & 求解/\allowbreak{}问题一/\allowbreak{}\mbox{问题一.py} & 附件 A & 质量评分、回归系数、推荐配比\\
二 & 求解/\allowbreak{}问题二/\allowbreak{}\mbox{问题二.py} & 附件 B、问题一结果 & 标度律参数、验证指标、配比修正\\
三 & 求解/\allowbreak{}问题三/\allowbreak{}\mbox{问题三.py} & 前两问共享结果、附件 C7 & 三档配置、预算扫描与敏感性\\
四 & 求解/\allowbreak{}问题四/\allowbreak{}\mbox{问题四.py} & 附件 C & 任务聚合、贡献分解与前沿预测\\
\bottomrule
\end{tabularx}
\end{table}
\FloatBarrier

\Needspace{8\baselineskip}
\subsubsection*{问题一：质量评价、冲突统计与配比建模}
\lstinputlisting[language=Python,firstnumber=1,caption={问题一完整求解程序},label={code:p1}]{../求解/问题一/问题一.py}

\Needspace{8\baselineskip}
\subsubsection*{问题二：广义标度律、验证与边际分析}
\lstinputlisting[language=Python,firstnumber=1,caption={问题二完整求解程序},label={code:p2}]{../求解/问题二/问题二.py}

\Needspace{8\baselineskip}
\subsubsection*{问题三：预算约束优化与敏感性分析}
\lstinputlisting[language=Python,firstnumber=1,caption={问题三完整求解程序},label={code:p3}]{../求解/问题三/问题三.py}

\Needspace{8\baselineskip}
\subsubsection*{问题四：任务解析、能力分解与前沿预测}
\lstinputlisting[language=Python,firstnumber=1,caption={问题四完整求解程序},label={code:p4}]{../求解/问题四/问题四.py}
